\chapter{As fases da Lua}\label{ch:g7:moon-phases}

Anos atrás você manteve um diário da \omterm{def:g2:moon-in-the-sky:moon}{Lua} e viu a forma prateada
marchar da \omterm{ex:g2:moon-in-the-sky:shapes}{lua crescente} à cheia e de volta, com o ``porquê''
prometido para depois. Depois é agora. Você tem tudo o que a
explicação precisa: uma \omterm{def:g2:moon-in-the-sky:moon}{Lua} que orbita todo mês, \omterm{def:g1:light-and-shadows:source}{luz} do Sol em linhas
retas e uma esfera que está sempre --- sempre --- iluminada pela
metade.

\section{O único fato por baixo de tudo}

\begin{proposition}[Metade iluminada, sempre]\label{prop:g7:moon-phases:halflit}
O Sol ilumina uma metade da \omterm{def:g2:moon-in-the-sky:moon}{Lua} a todo instante --- a metade voltada
para ele --- exatamente como ilumina metade da Terra. A \omterm{def:g2:moon-in-the-sky:moon}{Lua} não muda
de forma, não cresce nem encolhe: ela é uma bola de pedra vestindo um
lado de dia permanente e um lado de noite permanente, com a fronteira
deslizando conforme ela orbita. Tudo o que chamamos de ``\omterm{def:g7:moon-phases:phases}{fase}'' é
apenas a \emph{vista}: quanto do lado de dia calha de estar voltado
para a Terra.
\end{proposition}

\begin{definition}[Fases]\label{def:g7:moon-phases:phases}
As \emph{fases}\index{fase} da \omterm{def:g2:moon-in-the-sky:moon}{Lua} são as formas aparentes dela ao
longo do mês: \emph{lua nova} (lado de dia inteiramente virado para
longe --- não vemos nada), \emph{\omterm{ex:g2:moon-in-the-sky:shapes}{lua crescente}}, \emph{quarto
crescente} (metade do disco iluminada), \emph{gibosa} (mais que
metade), \emph{\omterm{ex:g2:moon-in-the-sky:shapes}{lua cheia}} (lado de dia inteiramente voltado para
\omterm{def:g5:series-parallel-circuits:parallel}{nós}), depois gibosa de novo, \emph{quarto minguante} e novamente uma
foice fina, fechando o ciclo em cerca de $29.5$ dias. Da nova à cheia
a parte iluminada cresce noite após noite: a \omterm{def:g2:moon-in-the-sky:moon}{Lua} está
\emph{crescente}\index{crescente}; da cheia de volta à nova ela
encolhe: \emph{minguante}\index{minguante}.
\end{definition}

\begin{omfigure}
\begin{tikzpicture}[scale=0.8]
  % dashed orbit
  \draw[dashed, omExo!70] (0,0) ellipse (4.2 and 3.1);
  % eight stations: dark half toward the left, sunlit half (white)
  % toward the sunlight arriving from the right
  \foreach \a in {0,45,...,315} {
    \begin{scope}[shift={({4.2*cos(\a)},{3.1*sin(\a)})}]
      \fill[omExo] (0,0) circle (0.46);
      \begin{scope}
        \clip (0,0) circle (0.46);
        \fill[white] (0,-0.5) rectangle (0.5,0.5);
      \end{scope}
      \draw[thick] (0,0) circle (0.46);
      \draw[thick] (0,0.46) -- (0,-0.46);
    \end{scope}
  }
  % the Earth, its own day side facing the sunlight
  \fill[omDef!35] (0,0) circle (0.55);
  \begin{scope}
    \clip (0,0) circle (0.55);
    \fill[omDef!12] (0,-0.6) rectangle (0.6,0.6);
  \end{scope}
  \draw[thick] (0,0) circle (0.55);
  \node[font=\footnotesize] at (0,0) {Terra};
  \foreach \y in {-2.4,-1.2,0,1.2,2.4} {
    \draw[-Stealth, omThm, thick] (7.7,\y) -- (6.4,\y);
  }
  \node[right, font=\small, omThm] at (7.8,0.55) {luz do Sol};
  \node[left, font=\footnotesize] at (-4.85,0) {cheia};
  \node[right, font=\footnotesize] at (4.85,0) {nova};
  \node[above, font=\footnotesize] at (0,3.75) {quarto minguante};
  \node[below, font=\footnotesize] at (0,-3.75) {quarto crescente};
\end{tikzpicture}

{\small O segredo inteiro numa figura: oito \omterm{def:g3:sun-days-seasons:year}{estações} da \omterm{def:g2:moon-in-the-sky:moon}{Lua} em
\omterm{def:g6:sun-earth-moon:axisorbit}{órbita}, cada uma iluminada pela metade do lado do Sol. O que muda é
apenas qual fatia do lado de dia fica voltada para a Terra.}
\end{omfigure}

\begin{method}[Seja a Terra, segure a Lua]\label{met:g7:moon-phases:walk}
A explicação, encenada em um minuto:
\begin{enumerate}
  \item num cômodo escuro, uma única lâmpada nua é o Sol; você é a
        Terra; uma bola branca na sua mão estendida é a \omterm{def:g2:moon-in-the-sky:moon}{Lua};
  \item fique de frente para a lâmpada, com a bola voltada para ela:
        a metade iluminada da bola fica virada para longe de você ---
        você vê o lado escuro dela: lua nova;
  \item gire devagar sem sair do lugar (a bola orbitando você): uma
        fatia iluminada aparece na borda da bola --- \omterm{ex:g2:moon-in-the-sky:shapes}{lua crescente}
        fininha; depois de um quarto de volta, metade da bola aparece
        iluminada: quarto \omterm{def:g7:moon-phases:phases}{crescente};
  \item continue girando: gibosa --- e depois, com a lâmpada nas suas
        costas e a bola inteiramente iluminada à sua frente: \omterm{ex:g2:moon-in-the-sky:shapes}{lua
        cheia}; complete a volta pela metade \omterm{def:g7:moon-phases:phases}{minguante} e pela foice
        fina até a nova.
\end{enumerate}
Uma volta sua reencena um mês de céu.
\end{method}

\section{A ideia errada, aposentada}

\begin{proposition}[As fases não são a sombra da Terra]\label{prop:g7:moon-phases:notshadow}
A resposta errada mais popular da astronomia: ``as \omterm{def:g7:moon-phases:phases}{fases} são a \omterm{def:g1:light-and-shadows:shadow}{sombra}
da Terra sobre a \omterm{def:g2:moon-in-the-sky:moon}{Lua}''. A geometria a refuta de três jeitos. A \omterm{def:g1:light-and-shadows:shadow}{sombra}
da Terra aponta \emph{para longe} do Sol --- ela só pode tocar uma
\omterm{def:g2:moon-in-the-sky:moon}{Lua} que esteja do lado oposto ao Sol, o que acontece exclusivamente
na \omterm{ex:g2:moon-in-the-sky:shapes}{lua cheia} (e mesmo aí só nos raros eclipses). Uma \omterm{def:g2:moon-in-the-sky:moon}{Lua} em foice
fica \emph{do lado do Sol} em relação a \omterm{def:g5:series-parallel-circuits:parallel}{nós}, longe demais da nossa
\omterm{def:g1:light-and-shadows:shadow}{sombra}. E a borda da \omterm{def:g1:light-and-shadows:shadow}{sombra} sobre a \omterm{def:g2:moon-in-the-sky:moon}{Lua} é sempre parte de um círculo
grande --- ela jamais poderia esculpir a fronteira de curvatura dupla
da forma gibosa. As \omterm{def:g7:moon-phases:phases}{fases} não precisam de \omterm{def:g1:light-and-shadows:shadow}{sombra} nenhuma: uma bola
iluminada pela metade, vista de lados diferentes, dá conta de tudo.
\end{proposition}

\begin{example}[Eclipse contra fase]\label{ex:g7:moon-phases:eclipsevs}
Mantenha os dois escurecimentos da \omterm{def:g2:moon-in-the-sky:moon}{Lua} em gavetas separadas.
\emph{\omterm{def:g7:moon-phases:phases}{Fase}}: o ciclo mensal lento, o próprio lado noturno da \omterm{def:g2:moon-in-the-sky:moon}{Lua} se
virando para \omterm{def:g5:series-parallel-circuits:parallel}{nós} --- nenhuma \omterm{def:g1:light-and-shadows:shadow}{sombra} nossa envolvida. \emph{\omterm{def:g7:shadows-eclipses:lunar}{Eclipse
lunar}}: uma hora rara em que a \omterm{ex:g2:moon-in-the-sky:shapes}{lua cheia} atravessa a nossa \omterm{prop:g7:shadows-eclipses:umbra}{umbra} ---
a \omterm{def:g1:light-and-shadows:shadow}{sombra} da Terra caindo de verdade sobre ela, cor de cobre e rápida.
O ritmo do diário é uma vista; o eclipse é uma visita.
\end{example}

\section{Ler o céu como um relógio}

\begin{example}[Onde cada fase guarda o horário dela]\label{ex:g7:moon-phases:hours}
O diagrama prende cada \omterm{def:g7:moon-phases:phases}{fase} a uma posição --- e cada posição a um
horário. A \omterm{ex:g2:moon-in-the-sky:shapes}{lua cheia} fica do lado oposto ao Sol: ela precisa nascer
no \omterm{def:g1:day-and-night:daynight}{pôr do sol}, reinar a noite inteira e se pôr no \omterm{def:g1:day-and-night:daynight}{nascer do sol}. A
lua nova viaja \emph{com} o Sol: fica no céu o dia inteiro, perdida
no brilho --- e é por isso que ninguém ``vê'' uma lua nova. O quarto
\omterm{def:g7:moon-phases:phases}{crescente} fica um quarto de volta atrás do Sol: nasce perto do
meio-dia e se põe perto da meia-noite --- a \omterm{def:g2:moon-in-the-sky:moon}{Lua} de depois do jantar.
A velha surpresa da \omterm{def:g2:moon-in-the-sky:moon}{Lua} de dia, da sua infância, agora é um horário
de trabalho: qualquer \omterm{def:g2:moon-in-the-sky:moon}{Lua} que não seja nova nem cheia divide parte do
turno dela com o Sol.
\end{example}

\begin{example}[Para que lado aponta a foice?]\label{ex:g7:moon-phases:points}
O lado iluminado de qualquer \omterm{def:g7:moon-phases:phases}{fase} fica voltado para o Sol --- a foice
é um arco apontado para o arqueiro dela. Foice do começo da noite, a
oeste: as pontas apontando para cima e para longe do clarão do \omterm{def:g1:day-and-night:daynight}{pôr do
sol}, abaixo do horizonte. O céu nunca desenha errado; os pintores,
muitas vezes, sim. Uma foice com as pontas voltadas para o Sol é o
caça-erro predileto da astronomia.
\end{example}

\begin{remark}[O rosto que nunca se vira]\label{rem:g7:moon-phases:sameface}
Mais uma observação do diário, enfim homenageada: as manchas escuras
da \omterm{def:g2:moon-in-the-sky:moon}{Lua} formam o mesmo rosto todas as noites, \omterm{def:g7:moon-phases:phases}{fase} após \omterm{def:g7:moon-phases:phases}{fase}, ano após
ano. A \omterm{def:g2:moon-in-the-sky:moon}{Lua} gira exatamente uma vez em torno do \omterm{def:g6:sun-earth-moon:axisorbit}{eixo} dela a cada
\omterm{def:g6:sun-earth-moon:axisorbit}{órbita} --- então a mesma metade fica voltada para a Terra para
sempre. Ninguém viu o lado oculto até uma sonda espacial dar a volta
com uma câmera, dentro do tempo de vida dos seus avós. Por que o giro
e a \omterm{def:g6:sun-earth-moon:axisorbit}{órbita} combinam tão perfeitamente é uma história profunda sobre
marés, guardada para os anos de graduação deste curso.
\end{remark}

\section{Exercícios}

\begin{exercise}[$\star$]\label{exo:g7:moon-phases:1}
Que fração da \omterm{def:g2:moon-in-the-sky:moon}{Lua} está iluminada pelo Sol a cada instante? O que é,
então, uma ``\omterm{def:g7:moon-phases:phases}{fase}''?
\end{exercise}

\begin{exercise}[$\star$]\label{exo:g7:moon-phases:2}
Recite as \omterm{def:g7:moon-phases:phases}{fases} em ordem a partir da lua nova e dê a \omterm{def:g2:measuring-time:duration}{duração} do
ciclo. Qual metade do ciclo é a ``\omterm{def:g7:moon-phases:phases}{crescente}''?
\end{exercise}

\begin{exercise}[$\star$]\label{exo:g7:moon-phases:3}
No método da caminhada orbital, onde ficam a lâmpada, a bola e você
na lua nova? E na \omterm{ex:g2:moon-in-the-sky:shapes}{lua cheia}?
\end{exercise}

\begin{exercise}[$\star$]\label{exo:g7:moon-phases:4}
Dê duas das três refutações geométricas de ``as \omterm{def:g7:moon-phases:phases}{fases} são a \omterm{def:g1:light-and-shadows:shadow}{sombra} da
Terra''.
\end{exercise}

\begin{exercise}[$\star$]\label{exo:g7:moon-phases:5}
Qual é a diferença entre a escuridão da lua nova e a escuridão de um
\omterm{def:g7:shadows-eclipses:lunar}{eclipse lunar}?
\end{exercise}

\begin{exercise}[$\star$]\label{exo:g7:moon-phases:6}
Por que a \omterm{ex:g2:moon-in-the-sky:shapes}{lua cheia} sempre nasce perto do \omterm{def:g1:day-and-night:daynight}{pôr do sol}? Argumente pela
posição dela no diagrama.
\end{exercise}

\begin{exercise}[$\star$]\label{exo:g7:moon-phases:7}
Por que ninguém jamais ``vê a lua nova nascer à meia-noite''? Onde
está a lua nova à meia-noite?
\end{exercise}

\begin{exercise}[$\star$]\label{exo:g7:moon-phases:8}
Uma foice do começo da noite fica pendurada a oeste. Para que lado
apontam as pontas dela em relação ao Sol que sumiu --- e por que uma
foice nunca pode apontar as pontas para o Sol?
\end{exercise}

\begin{exercise}[$\star\star$]\label{exo:g7:moon-phases:9}
Hoje à noite a \omterm{def:g2:moon-in-the-sky:moon}{Lua} está gibosa \omterm{def:g7:moon-phases:phases}{crescente}. Onde ela está na \omterm{def:g6:sun-earth-moon:axisorbit}{órbita}
dela? Mais ou menos quando ela nasceu, e qual \omterm{def:g7:moon-phases:phases}{fase} --- e qual horário
--- chega cinco noites depois?
\end{exercise}

\begin{exercise}[$\star\star$]\label{exo:g7:moon-phases:10}
Um \omterm{def:g7:shadows-eclipses:solar}{eclipse solar} só pode acontecer numa \omterm{def:g7:moon-phases:phases}{fase}, e um \omterm{def:g7:shadows-eclipses:lunar}{eclipse lunar} só
em outra. Diga as duas, justifique com a geometria do alinhamento e
explique por que, mesmo assim, os eclipses são raros nessas \omterm{def:g7:moon-phases:phases}{fases}.
\end{exercise}

\begin{exercise}[$\star\star$]\label{exo:g7:moon-phases:11}
Astronautas no lado visível da \omterm{def:g2:moon-in-the-sky:moon}{Lua} olham para cima, para a Terra.
Descreva as ``\omterm{def:g7:moon-phases:phases}{fases} da Terra'' que eles veem ao longo de um mês e qual
\omterm{def:g7:moon-phases:phases}{fase} a Terra mostra a eles quando vemos uma lua nova. (Rode a
geometria nos dois sentidos.)
\end{exercise}

\begin{exercise}[$\star\star\star$]\label{exo:g7:moon-phases:12}
O segredo da foice antiga: em noites limpas, a parte \emph{escura} da
\omterm{def:g2:moon-in-the-sky:moon}{Lua} em foice brilha de leve --- ``a lua velha nos braços da lua
nova''. Você conheceu isso como charada no segundo ano; resolva agora
por completo: percorra o caminho da \omterm{def:g1:light-and-shadows:source}{luz}, explique por que a \omterm{def:g1:light-and-shadows:source}{luz}
cinzenta é mais forte exatamente nas \omterm{def:g7:moon-phases:phases}{fases} de foice (qual \omterm{def:g7:moon-phases:phases}{fase} a
\emph{Terra} mostra à \omterm{def:g2:moon-in-the-sky:moon}{Lua} nesses momentos?) e por que ela some na
direção da \omterm{ex:g2:moon-in-the-sky:shapes}{lua cheia}.
\end{exercise}

\section{Problema: Um mês com a Lua}

\begin{problem}[{Problema de fim de semana --- o mês aberto do
observatório; trinta noites de Lua, um cartaz de parede; o
questionário do monitor}]\label{pb:g7:moon-phases:1}
O observatório da cidade promove um ``Mês com a \omterm{def:g2:moon-in-the-sky:moon}{Lua}'' aberto ao
público. Como monitor júnior, você cuida do cartaz de parede e
responde aos visitantes.

\medskip\noindent\textbf{Parte I --- Montando o cartaz.}
\begin{enumerate}
  \item A noite 0 é a lua nova. Preencha os marcos do cartaz: em
        quais noites, mais ou menos, caem o quarto \omterm{def:g7:moon-phases:phases}{crescente}, a \omterm{ex:g2:moon-in-the-sky:shapes}{lua
        cheia} e o quarto \omterm{def:g7:moon-phases:phases}{minguante}?
  \item Ao lado de cada marco, o cartaz quer os horários de nascer.
        Associe: nasce no \omterm{def:g1:day-and-night:daynight}{pôr do sol}; \ nasce perto do meio-dia; \
        nasce perto da meia-noite; \ nasce e se põe com o Sol.
  \item Um visitante pergunta por que o cartaz mostra a foice baixa
        no céu do \emph{oeste} no começo da noite, mas a \omterm{ex:g2:moon-in-the-sky:shapes}{lua cheia}
        subindo pelo \emph{leste} ao anoitecer. Explique com as
        posições na \omterm{def:g6:sun-earth-moon:axisorbit}{órbita}.
  \item Outro visitante jura ter visto a \omterm{def:g2:moon-in-the-sky:moon}{Lua} às dez da manhã. Quais
        \omterm{def:g7:moon-phases:phases}{fases} do cartaz tornam possível uma \omterm{def:g2:moon-in-the-sky:moon}{Lua} de manhã?
\end{enumerate}

\medskip\noindent\textbf{Parte II --- As demonstrações do monitor.}
\begin{enumerate}[resume]
  \item Você faz a caminhada da lâmpada e da bola para um grupo
        escolar. Um aluno \omterm{def:g2:ice-water-steam:meltfreeze}{congela} você no meio da volta: a bola
        aparece três quartos iluminada. Diga o nome da \omterm{def:g7:moon-phases:phases}{fase} e
        localize a lâmpada, a Terra-você e a bola para o grupo.
  \item ``Então a \omterm{def:g1:light-and-shadows:shadow}{sombra} da Terra é que faz as \omterm{def:g7:moon-phases:phases}{fases}?'', pergunta o
        aluno, apontando para o lado escuro da sua bola. Dê a
        refutação da própria demonstração: de quem é o lado noturno
        que eles estão vendo, e onde a \omterm{def:g1:light-and-shadows:shadow}{sombra} da Terra realmente
        estaria?
  \item O grupo pergunta por que sempre veem ``o coelho'' (ou ``o
        rosto'') na \omterm{def:g2:moon-in-the-sky:moon}{Lua}, e nunca as costas dela. Responda com o fato
        do mesmo rosto.
  \item Um pai cético: ``Se a \omterm{def:g2:moon-in-the-sky:moon}{Lua} gira, como você afirma, não
        deveríamos vê-la girando?'' Resolva a contradição aparente:
        como a \omterm{def:g2:moon-in-the-sky:moon}{Lua} consegue girar exatamente uma vez por \omterm{def:g6:sun-earth-moon:axisorbit}{órbita} e,
        com isso, \emph{esconder} de \omterm{def:g5:series-parallel-circuits:parallel}{nós} o giro dela?
\end{enumerate}

\medskip\noindent\textbf{Parte III --- Noite de questionário.}
\begin{enumerate}[resume]
  \item Carta A: ``A \omterm{def:g2:moon-in-the-sky:moon}{Lua} de hoje nasceu à meia-noite e mostra a
        metade \emph{esquerda} iluminada (vista rumo à manhã).
        \omterm{def:g7:moon-phases:phases}{Crescente} ou \omterm{def:g7:moon-phases:phases}{minguante} --- e qual quarto?''
  \item Carta B: ``Um cartaz de artista mostra uma \omterm{def:g2:moon-in-the-sky:moon}{Lua} em foice à
        meia-noite, alta no céu, com as pontas apontando direto para
        baixo, na direção das cores do \omterm{def:g1:day-and-night:daynight}{pôr do sol} no horizonte.''
        Liste todos os erros astronômicos do cartaz.
  \item Carta C: ``Durante a \omterm{ex:g2:moon-in-the-sky:shapes}{lua cheia} de hoje, os visitantes também
        poderiam presenciar um \omterm{def:g7:shadows-eclipses:solar}{eclipse solar} em algum lugar da
        Terra?'' Resolva com a geometria das \omterm{def:g7:moon-phases:phases}{fases}.
  \item Discurso de encerramento: em três frases, dê o único fato de
        base do mês, a ideia errada que o mês enterrou e o truque de
        horário que os visitantes devem levar para casa (o que
        qualquer \omterm{def:g7:moon-phases:phases}{fase} diz sobre a hora em que ela nasce).
\end{enumerate}
\end{problem}
